\subsection{Polynomial Functions} 
\begin{definition}[Polynomial Function]
A function is a polynomial function if its expression is a polynomial.
\end{definition}
\begin{example}
Let $f(x)=2x^3-5x^2+4x+7$. Evaluate:
\begin{enumerate}[a.]
\item \makeblankfill{$f(3)={}$ }
\item \makeblankfill{$f(0)={}$ }
\item \makeblankfill{$f(-2)={}$ }
\end{enumerate}
\end{example}
Polynomials give an insight into the Order of Operations.
\begin{definition}[Order of Operations]
In expressions with more than one operation, they should be completed in the following order:
\begin{enumerate*}
\item Exponents
\item Multiplication (and division)
\item Addition (and subtraction)
\end{enumerate*}
To violate this order, use grouping symbols (like parentheses).
\end{definition}
The reason that this order is chosen is to make polynomials work as expected. To evaluate $f(-2)$, we:
\begin{enumerate*}
\item found each power of $-2$
\item multiplied by its coefficient
\item then added up all the terms
\end{enumerate*}
This is why \makeblank[1.25in] come before \makeblank[1.25in] in order of operations: we want $-x^2$ to mean \makeblank[1.25in] the number $x$, then \makeblank[1.25in].
\begin{example}
Evaluate each of the following.\lbr
$2^2=\makeblank[1in]$\hfill
$-2^2=\makeblank[1in]$\hfill
$(-2)^2=\makeblank[1in]$\hfill
\end{example}